\documentclass[10pt,a4paper]{amsart}
\usepackage[margin=19mm]{geometry}
\usepackage{amsmath,amssymb,mathtools}
\usepackage[scr=rsfs,cal=euler]{mathalfa}
\usepackage{xfrac}
\usepackage[unicode,hidelinks,bookmarks=false]{hyperref}
\newcommand{\ie}{i.e.\ }
\newcommand{\cf}{cf.\ }
\newcommand{\resp}{respectively\ }
\newcommand{\Subsection}{\subsection}
\providecommand{\nobreakdash}{\nobreak\hbox{-}}
\setlength{\emergencystretch}{2em}
\multlinegap0pt
\usepackage{bm}
\usepackage{bbm}
\usepackage{dsfont}
\usepackage{xspace}
\usepackage{cancel}
\usepackage{mathtools}
\usepackage{amsthm}
\makeatletter
\@ifundefined{thm}{\newtheorem{thm}{Theorem}}{}
\@ifundefined{lem}{\newtheorem{lem}[thm]{Lemma}}{}
\makeatother
\newcommand{\mbb}{\mathbb}
\newcommand{\mc}{\mathcal}
\title[Pareto-Optimal Linear Programming]{Pareto-Optimal Linear Programming}
\author{Bart van Rossum, Twan Dollevoet}
\date{Source: arXiv:2509.18073v1 (CC BY 4.0)}
\begin{document}
\maketitle

\noindent\emph{Source and licence.} Pareto-Optimal Linear Programming. Version arXiv:2509.18073v1 (CC BY 4.0), distributed under the Creative Commons Attribution 4.0 International licence, as recorded in the arXiv licence metadata for this exact version and shown on the abstract page https://arxiv.org/abs/2509.18073v1. Full rights notice: licenses/arxiv-2509.18073-CC-BY-4.0.txt.\par\smallskip

\noindent\textit{Editorial scope.} Counted unit: the Thm whose source label is \texttt{None} in the article (source0.tex lines 350--352, source label \texttt{None}), delivered together with the article's own complete original proof of it (source0.tex lines 354--379, one \texttt{proof} environment of 26 lines). No mathematics is written, repaired or re-derived: statement and proof are the article's own text line for line. Required context retained uncounted: lem:opt\_vertex (lines 215--218); thm:fpo\_matching (lines 328--331). Every definition, display and label of those lines is retained unchanged. Omitted from this package and not redistributed: 1--214, 219--327, 332--349, 353--353, 380--500. Nothing inside a retained excerpt is omitted or altered: every retained body is delivered exactly as the article's lines are written, so no omission receipt of any kind accompanies this package. Citations: the retained excerpts carry no citation command at all, so no bibliography entry of the article is transcribed and no reference is redistributed. Reference adaptation: none was needed and none was made; every reference inside the delivered unit resolves inside it, so no unresolved reference or citation is produced. Printed slips kept literal: no printed word, symbol, hypothesis or number of the counted statement or of its proof was altered. Layout and class: the article's own class is \texttt{article} with packages including geometry, parskip, appendix, array, bm, caption, enumerate, rotating, hyperref, blkarray, bigstrut, multirow, booktabs, siunitx; the wrapper is the lane's pinned \texttt{amsart} editorial wrapper at 10pt on A4, which restates the theorem-like environments the retained text needs and adds the article's own macro definitions that the retained text needs (\texttt{\textbackslash mbb}, \texttt{\textbackslash mc}), with extra wrapper packages (bm, bbm, dsfont, xspace, cancel, mathtools). The delivered numbering is the unit's own. Assets: no retained span contains a figure environment, a graphics-inclusion command, a PDF-page inclusion, a table or a TikZ picture; no image file is copied into this package, the package declares no asset dependency and no image file is redistributed with it. Licence: version arXiv:2509.18073v1 of this article is distributed under the Creative Commons Attribution 4.0 International licence, as recorded in the arXiv licence metadata for this exact version and shown on the abstract page https://arxiv.org/abs/2509.18073v1. A full rights notice is delivered at licenses/arxiv-2509.18073-CC-BY-4.0.txt. Characters: every retained excerpt is pure ASCII, so the delivered engine, XeLaTeX with its default Latin Modern text font, reports no missing character.
\begin{lem}
\label{lem:opt_vertex}
There always exists an optimal solution to \textsc{Max-Pareto} at a vertex of $\mc{X}$. 
\end{lem}

\begin{thm}
If matching $M$ in $G$ is PO, then $M$ is fPO. 
\label{thm:fpo_matching}
\end{thm}

\begin{thm}
The decision version of \textsc{Max-Pareto} is $\mc{NP}$-complete.
\end{thm}

\begin{proof}
It is readily seen that \textsc{Max-Pareto} is in $\mc{NP}$, as membership of $\bm{x} \in \mathcal{X}$ can be easily confirmed, and Pareto-optimality can be checked in polynomial time by solving the following linear program:
\begin{subequations}
\begin{align}
\max \quad & \bm{e}^\top U \bm{y} \\
\text{s.t.} \quad & U \bm{y} \geq U \bm{x} \\ 
& \bm{y} \in \mc{X}.  
\end{align}
\end{subequations}
We now reduce \textsc{Constrained Pareto-Optimal Matching} to an instance of the decision version of \textsc{Max-Pareto}: Does there exist an $\bm{x} \in \mc{X}_{P}$ for which $\bm{c}^\top \bm{x} \geq k$? We construct this instance of $\textsc{Max-Pareto}$ as follows.
\begin{itemize}
    \item \textit{Feasible region:} Construct a bipartite graph $G := (\mc{A} \cup \mc{O}, E, \bm{w})$ as follows. Edge $e = (i, a)$ is included in $E$ if agent $i \in \mc{A}$ admits item $a \in \mc{O}$, and has weight $w_{(i, a)} := 1 + | \{ o \in \mc{O} : a \succ_{P_i} o\} |$. In other words, we choose weights that assign more value to objects higher up in the agent's preference profile. The bipartite graph is not necessarily complete. We now define $\mc{X} := \{ \bm{x} \in \mathbb{R}^{|E|}_{+} : \sum\limits_{e \in \delta(v)} x_e \leq 1, \forall v \in V_1 \cup V_2\}$ as the matching polytope of $G$, where $\delta(v)$ is the set of edges incident to $v$. Since $G$ is bipartite, the corresponding constraint matrix is totally unimodular and each vertex of the matching polytope corresponds to an integral matching.
    \item \textit{Payoff mapping:} There are $|\mc{A}|$ agents. For any $e = (i, a) \in E$, we set $U_{(i, e)} := w_e$. This payoff function ensures that an integral matching is Pareto-optimal in \textsc{Max-Pareto} if and only if it is Pareto-optimal in \textsc{Constrained Pareto-Optimal Matching}.
    \item \textit{Objective function:} Choose $\bm{c} \in \mbb{R}^{|E|}$ by setting $c_{(i, a)} := 1$ if $a \in Q$ and $c_{(i, a)} := 0$ otherwise. The objective is thus to maximise the number of matched objects in $Q$.
    \item \textit{Objective threshold:} The objective threshold equals $k := |Q|$.
\end{itemize} 
This instance of \textsc{Max-Pareto} is constructed in time polynomial in the size of the input. We now show that it is a yes-instance if and only if \textsc{Constrained Pareto-Optimal Matching} is a yes-instance.

$\Rightarrow$ First, suppose that \textsc{Constrained Pareto-Optimal Matching} is a yes-instance. Let $M$ be a Pareto-optimal matching containing all items in $Q$. Choose $\bm{x}$ by setting $x_{e} := 1$ if $e \in M$, and $x_e := 0$ otherwise. It is clear that $\bm{x}$ is feasible in $\mc{X}$ and meets the objective threshold. 

It remains to show that $\bm{x}$ satisfies Pareto-optimality in \textsc{Max-Pareto}. By construction, this is equivalent to showing that $M$ is fPO with respect to all matchings in $G$. This follows directly from Theorem~\ref{thm:fpo_matching} and the fact that $M$ is PO. We conclude that \textsc{Max-Pareto} is a yes-instance.

$\Leftarrow$ Now, suppose that \textsc{Max-Pareto} is a yes-instance with corresponding solution $\bm{x}$. We aim to convert this to a Pareto-optimal matching $M$ that matches all items in $Q$. By Lemma~\ref{lem:opt_vertex}, we know there exists an optimal vertex solution $\bm{x}'$. By total unimodularity of the bipartite matching polytope, this solution corresponds to an integral matching. This matching is not Pareto-dominated by any point in $\mc{X}$. In particular, it is not dominated by any integral matching. Since $\bm{c}^\top \bm{x} \geq k$, the matching contains all objects in $Q$. Hence, \textsc{Constrained Pareto-Optimal Matching} is a yes-instance.  

Since \textsc{Constrained Pareto-Optimal Matching} is $\mc{NP}$-hard and can be reduced to \textsc{Max-Pareto} in polynomial time, we conclude that \textsc{Max-Pareto} is $\mc{NP}$-hard as well.
\end{proof}

\end{document}
